\subsection{Đánh Giá Mô Hình (Evaluation)}


\subsubsection{So sánh mô hình tuyến tính}

\begin{table}[H]
\centering
\caption{So sánh các mô hình trên tập test}
\label{tab:part2_model_summary}
\small
\renewcommand{\arraystretch}{1.0}
\begin{adjustbox}{max width=\textwidth}
\begin{tabular}{@{}lrrrrr@{}}
\hline
\textbf{Mô hình} & \textbf{Test $R^2$} & \textbf{Test RMSE} & \textbf{Test MAE} &
\textbf{CV-MSE} & \textbf{CV-$R^2$} \\
\hline
Lasso & 0.748403 & 0.056872 & 0.037895 & 0.005074 & 0.673462 \\
OLS full & 0.747794 & 0.056941 & 0.037952 & 0.005075 & 0.673407 \\
Ridge & 0.747678 & 0.056954 & 0.038204 & 0.005074 & 0.673433 \\
OLS selected & 0.743913 & 0.057378 & 0.039013 & 0.005103 & 0.671611 \\
\hline
\end{tabular}
\end{adjustbox}
\end{table}

\begin{figure}[H]
    \centering
    \includegraphics[width=\textwidth]{assets/part2/model_comparison.png}
    \caption{So sánh Test $R^2$, RMSE và MAE giữa các mô hình tuyến tính. Lasso đạt Test $R^2$
    cao nhất, nhưng khác biệt giữa Lasso, OLS full và Ridge rất nhỏ.}
    \label{fig:part2_model_comparison}
\end{figure}

Nhận xét:
\begin{itemize}
    \item Lasso có Test $R^2$ cao nhất là 0.748403 và Test RMSE thấp nhất là 0.056872.
    \item OLS full gần như tương đương Lasso, cho thấy quan hệ tuyến tính sau tiền xử lý đã mô tả
          dữ liệu khá tốt.
    \item Ridge có CV-MSE thấp nhất trong nhóm mô hình tuyến tính, phản ánh regularization giúp
          ổn định hơn trên các fold train.
    \item OLS selected ít biến hơn nhưng Test $R^2$ giảm nhẹ xuống 0.743913; việc rút gọn biến giúp
          dễ diễn giải hơn nhưng có thể mất một phần thông tin dự báo.
\end{itemize}

Sự khác biệt giữa Lasso, OLS full và Ridge rất nhỏ: Test $R^2$ của ba mô hình chỉ chênh nhau ở
chữ số thập phân thứ tư. Lý do chính là pipeline đã làm phần việc mà regularization thường hỗ trợ:
log-transform giảm lệch phải, Winsorization giảm ảnh hưởng cực trị, và VIF filtering đã loại
\texttt{log\_pl\_orbper}, \texttt{log\_pl\_eqt} để đưa VIF cuối cùng xuống dưới 10. Khi đa cộng
tuyến đã được xử lý trước, Ridge không còn lợi thế ổn định hệ số quá lớn, còn Lasso với
$\lambda^*=0.1$ vẫn không tạo nghiệm thưa. Vì vậy mô hình tuyến tính cơ bản đã gần đạt
hiệu năng tối đa của nhóm mô hình tuyến tính trên tập feature hiện có.

Mặt khác, $R^2_{\mathrm{test}}\approx0.748$ cũng có nghĩa khoảng 25.2\% phương sai của
\texttt{log\_pl\_rade} vẫn chưa được giải thích. Phần còn lại có thể đến từ ba nguồn: quan hệ vật
lý phi tuyến giữa khối lượng--bán kính ở từng cấu tạo hành tinh, thiếu biến mô tả thành phần bên
trong/tuổi hệ sao/độ chính xác phép đo, và selection bias của dữ liệu quan trắc thiên văn. Do đó,
kết quả này nên được đọc như một mô hình nội suy tuyến tính tốt cho nhóm terrestrial/Super-Earth,
không phải một phương trình vật lý đầy đủ của quá trình hình thành hành tinh.

\subsubsection{Phân tích hệ số và feature importance}

\begin{figure}[H]
    \centering
    \includegraphics[width=0.9\textwidth]{assets/part2/feature_importance.png}
    \caption{Feature importance dựa trên hệ số Ridge sau chuẩn hóa. Vì các feature đã được
    standardize, độ lớn hệ số có thể so sánh trực tiếp hơn.}
    \label{fig:part2_feature_importance}
\end{figure}

Các biến quan trọng nhất gồm \texttt{log\_pl\_bmasse}, \texttt{log\_sy\_dist} và
\texttt{log\_pl\_trandep}. Trong đó, \texttt{log\_pl\_bmasse} là biến có ảnh hưởng dương mạnh nhất,
phù hợp với quan hệ vật lý giữa khối lượng và bán kính hành tinh.
\texttt{log\_sy\_dist} cũng có hệ số dương khá lớn, nhưng không nên diễn giải theo nghĩa nhân quả
rằng hành tinh ở xa Trái Đất thì tự nhiên có bán kính lớn hơn. Cách đọc hợp lý hơn là
\emph{selection bias}: ở khoảng cách xa, các kính thiên văn dễ phát hiện và xác nhận những hành
tinh có tín hiệu mạnh hơn, thường là transit sâu hơn hoặc bán kính lớn hơn. Vì vậy hệ số này phản
ánh thiên lệch quan sát của catalog nhiều hơn là một cơ chế vật lý trực tiếp.

\subsubsection{Phân tích phần dư và giả định mô hình}

\begin{figure}[H]
    \centering
    \includegraphics[width=\textwidth]{assets/part2/residual_plots.png}
    \caption{Bốn biểu đồ chẩn đoán phần dư của mô hình Lasso: Residuals vs Fitted, Normal Q-Q,
    Scale-Location và Cook's Distance.}
    \label{fig:part2_residual}
\end{figure}

Nhận xét từ residual plots:
\begin{itemize}
    \item Phần dư tập trung quanh 0, cho thấy mô hình không có thiên lệch hệ thống quá rõ trên đa
          số quan sát.
    \item Q-Q plot cho thấy phần giữa của phân phối phần dư tương đối gần chuẩn, nhưng hai đuôi vẫn
          lệch khỏi đường chéo; đây là dấu hiệu dữ liệu còn ngoại lai hoặc quan hệ phi tuyến.
    \item Cook's Distance lớn nhất được ghi nhận là 0.617979. Với $n=217$ mẫu test, ngưỡng
          $4/n \approx 0.01843$ đánh dấu 17 điểm influential.
\end{itemize}

\begin{table}[H]
\centering
\caption{Kiểm định chẩn đoán trên phần dư test của Lasso}
\label{tab:part2_residual_tests}
\begin{tabular}{@{}lrrp{0.45\textwidth}@{}}
\hline
\textbf{Kiểm định} & \textbf{Statistic} & \textbf{p-value} & \textbf{Diễn giải} \\
\hline
Shapiro-Wilk & 0.887620 & $1.21\times10^{-11}$ &
Bác bỏ giả thuyết phần dư phân phối chuẩn; hai đuôi còn lệch, phù hợp với Q-Q plot. \\
Breusch-Pagan & 36.895113 & 0.000120 &
Bác bỏ giả thuyết phương sai không đổi; vẫn còn heteroskedasticity sau log-transform và
Winsorization. \\
\hline
\end{tabular}
\end{table}

Nhóm không xóa 17 điểm influential vì đây là dữ liệu thiên văn thật, không phải lỗi nhập liệu đã
được chứng minh. Các điểm này có thể đại diện cho hệ hành tinh hiếm, sai số đo lớn, hoặc vùng biên
của catalog. Xóa cứng sẽ làm báo cáo đẹp hơn về mặt diagnostic nhưng dễ tạo thiên lệch mẫu. Vì
vậy nhóm giữ lại chúng, dùng log-transform/Winsorization để giảm ảnh hưởng cực trị, và ghi nhận
đây là một giới hạn khi diễn giải hệ số.

\subsubsection{Kiểm chứng giả thuyết Giãn nở nhiệt (Thermal Expansion Hypothesis)}

\paragraph{Kiểm định tác động tương tác giữa Khối lượng và Nhiệt độ:}
Từ hiện tượng đường LOWESS bị cong nhẹ ở vùng dự đoán cao trên biểu đồ chẩn đoán Residuals vs Fitted, 
giả thuyết được đặt ra là: tác động của khối lượng (lực nén hấp dẫn) lên cấu trúc bán kính hành tinh 
không độc lập, mà bị chi phối bởi mức độ bức xạ nhiệt (hiệu ứng giãn nở). Để kiểm chứng, mô hình hồi 
quy được mở rộng thành một cấu trúc phi tuyến cục bộ thông qua đại lượng tương tác:
\begin{equation}
    Y = \beta_1 X_{\text{mass}} + \beta_2 X_{\text{heat}} + \beta_3 (X_{\text{mass}} \times X_{\text{heat}})
\end{equation}

Thực nghiệm đầu tiên tiến hành bổ sung biến tương tác giữa khối lượng hành tinh và nhiệt độ bề mặt 
của sao chủ (\texttt{log\_pl\_bmasse} $\times$ \texttt{st\_teff}). Kết quả truy xuất cho thấy p-value 
của biến này đạt mức 0.098, vượt ngưỡng mức ý nghĩa $\alpha = 0.05$. Sự thất bại của giả thuyết này 
có thể lý giải bằng bản chất vật lý: nhiệt độ bề mặt sao chủ không phản ánh chính xác thông lượng 
nhiệt (heat flux) mà hành tinh thực nhận nếu khuyết đi tham số về khoảng cách quỹ đạo.

Để khắc phục, thực nghiệm thứ hai thay thế nhiệt độ sao bằng thông lượng bức xạ thực nhận (\texttt{log\_pl\_insol} 
tạo thành biến tương tác (\texttt{log\_pl\_bmasse} $\times$ \texttt{log\_pl\_insol}). Kết quả ghi 
nhận p-value $\approx 0$, khẳng định sự tồn tại của cơ chế tương tác nhiệt - khối lượng trên tập 
huấn luyện. Tuy nhiên, nghịch lý xuất hiện khi kiểm định hiệu năng: hệ số xác định trên tập kiểm thử 
($R^2_{\text{Test}}$) lại sụt giảm từ 0.747 xuống 0.741, đồng thời RMSE tăng.

Bản chất của hiện tượng này là minh chứng kinh điển cho bài toán \emph{Sự đánh đổi giữa Độ chệch và 
Phương sai (Bias-Variance Tradeoff)}. Biến tương tác đã vô tình kích hoạt hiện tượng Đa cộng tuyến 
cấu trúc (Structural Multicollinearity). Mặt phẳng OLS bị ép phải vặn xoắn (giảm Bias) nhằm cố gắng 
nội suy khít với các biến thiên cục bộ sinh ra từ hiệu ứng giãn nở nhiệt trong tập huấn luyện, từ đó 
đánh mất khả năng khái quát hóa quy luật (tăng Variance) khi tiếp xúc với dữ liệu kiểm thử.

\paragraph{Kết luận:}
Nỗ lực triệt tiêu sai số phần dư cục bộ bằng các biến tương tác đã dẫn đến hiện tượng quá khớp (Overfitting). 
Hệ quả này củng cố tính đúng đắn của nguyên lý tinh giản Occam's Razor trong thiết kế Machine Learning: 
cấu trúc OLS thuần túy hiện tại (không chứa đại lượng tương tác) là điểm cân bằng toàn cục tối ưu nhất 
giữa năng lực giải thích hiện tượng và sức mạnh ngoại suy định lượng.

\subsubsection{Phương pháp nâng cao: Kernel Ridge Regression}

Ngoài các mô hình tuyến tính chính, nhóm thử Kernel Ridge Regression với RBF kernel:
\begin{equation}
    k(\mathbf{x},\mathbf{x}')
    = \exp\left(-\gamma\|\mathbf{x}-\mathbf{x}'\|^2\right).
\end{equation}

Do Kernel Ridge cần ma trận Gram kích thước $n\times n$, chi phí huấn luyện tăng nhanh hơn
các mô hình tuyến tính. Với lần chạy cuối cùng trên Kaggle, nhóm dùng toàn bộ 867 mẫu train sau
tiền xử lý để fit Kernel Ridge. Hyperparameters được chọn trên validation tách từ train, sau đó
mới đánh giá một lần trên test.

\begin{table}[H]
\centering
\caption{Một số cấu hình Kernel Ridge theo validation RMSE}
\label{tab:part2_kernel_search}
\begin{tabular}{@{}rrr@{}}
\hline
\textbf{Alpha} & \textbf{Gamma} & \textbf{Validation RMSE} \\
\hline
0.010 & 0.010 & 0.037282 \\
0.001 & 0.001 & 0.037780 \\
0.100 & 0.010 & 0.038428 \\
0.001 & 0.010 & 0.040044 \\
0.010 & 0.050 & 0.043404 \\
0.010 & 0.001 & 0.044958 \\
0.100 & 0.050 & 0.045358 \\
1.000 & 0.010 & 0.046950 \\
\hline
\end{tabular}
\end{table}

\begin{table}[H]
\centering
\caption{Kết quả Kernel Ridge tốt nhất trên test set}
\label{tab:part2_kernel_result}
\begin{tabular}{@{}lrrrrr@{}}
\hline
\textbf{Mô hình} & \textbf{Alpha} & \textbf{Gamma} & \textbf{Test $R^2$} & \textbf{RMSE} & \textbf{MAE} \\
\hline
Kernel Ridge (RBF) & 0.01 & 0.01 & 0.855688 & 0.043072 & 0.029685 \\
\hline
\end{tabular}
\end{table}

So với mô hình tuyến tính tốt nhất là Lasso trong Bảng~\ref{tab:part2_model_summary},
Kernel Ridge cải thiện rõ rệt trên cùng test set: Test $R^2$ tăng từ 0.748403 lên 0.855688
(tăng 0.107285 điểm tuyệt đối), RMSE giảm từ 0.056872 xuống 0.043072, và MAE giảm từ
0.037895 xuống 0.029685. Nếu nhìn theo phần phương sai chưa giải thích được, Lasso còn
$1-0.748403=0.251597$, trong khi Kernel Ridge còn $1-0.855688=0.144312$; tức phần sai lệch
chưa giải thích giảm khoảng 42.6\%. Đây là mức chênh lệch lớn hơn nhiều so với khác biệt rất
nhỏ giữa OLS full, Ridge và Lasso, nên kết quả này không chỉ là dao động ngẫu nhiên giữa các
mô hình tuyến tính.

Nguyên nhân hợp lý là RBF kernel không ép dữ liệu nằm trên một mặt phẳng tuyến tính duy nhất.
Sau log-transform và VIF filtering, OLS/Ridge/Lasso đã xử lý tốt phần quan hệ gần tuyến tính,
nhưng các đặc trưng thiên văn như khối lượng hành tinh, chu kỳ quỹ đạo, nhiệt độ sao mẹ và khoảng
cách hệ sao vẫn có thể tương tác theo dạng cong cục bộ. Kernel Ridge đo độ giống nhau giữa các
hành tinh trong không gian feature đã chuẩn hóa, từ đó nội suy bán kính dựa trên các mẫu lân cận
phi tuyến. Việc cấu hình tốt nhất có $\gamma=0.01$ cũng cho thấy mô hình không chọn một kernel quá
hẹp để ghi nhớ từng điểm riêng lẻ, mà dùng một mức làm trơn vừa phải trên toàn bộ tập train.

Tuy nhiên, kết quả Kernel Ridge không làm mất vai trò của các mô hình tuyến tính chính. Thứ nhất,
OLS/Ridge/Lasso cho phép đọc trực tiếp hệ số, kiểm định giả thuyết, phân tích residuals, Cook's
Distance và selection bias; đây là các phần quan trọng của yêu cầu data fitting. Kernel Ridge cho
kết quả dự báo tốt hơn nhưng khó quy đổi thành ảnh hưởng riêng lẻ của từng biến. Thứ hai, mô hình
phải lưu ma trận Gram $n\times n$ và giải hệ tuyến tính kích thước $n$, làm chi phí bộ nhớ tăng
xấp xỉ $O(n^2)$ và chi phí giải hệ tăng xấp xỉ $O(n^3)$. Với 867 mẫu train hiện tại, việc chạy trên
Kaggle vẫn khả thi; nhưng nếu mở rộng về toàn bộ archive hoặc thêm nhiều đặc trưng hơn, chi phí này
sẽ tăng nhanh hơn nhiều so với nhóm mô hình tuyến tính. Vì vậy, nhóm xem Kernel Ridge là bằng chứng
bổ sung rằng dữ liệu có thành phần phi tuyến đáng kể, còn OLS/Ridge/Lasso vẫn là trục phân tích
chính vì minh bạch, nhẹ tài nguyên và phù hợp hơn cho phần diễn giải thống kê.

\FloatBarrier


